%% kapa: exam=2 round=19 year=2008 period=1 subject=appraisal_practice subject_ko=감정평가실무 seq=1 number=1 points=40
%% source: Q-Net 감정평가사 기출문제 (cst00309) artlSeq=5024177 file=extracted/1008/5024177/2008년 제19회 감정평가사 2차시험 기출문제/1교시.hwp
%% conversion: kordoc 4.24.51 -> tools/gen2.py (auto); review=proofread-sonnet; flags=tikz
\begin{longquestion}{【문제 1】}
감정평가사 L씨는 택지개발예정지구로 지정고시된 지역의 보상에 대하여 중앙토지수용위원회로부터 이의재결평가 의뢰를 받았다. 보상관련법규의 제규정 등을 참작하고 제시된 자료를 활용하여 보상액을 산정하시오. (40점) \par
\medskip\noindent (1) 의뢰토지에 대한 가격시점 결정 및 비교표준지 선정사유를 설명하고, 기호 3을 제외한 나머지 토지의 보상감정평가액을 산정하시오. \par
\medskip\noindent (2) 건물의 보상감정평가액을 산정하시오. \par
\medskip\noindent (3) <자료 8>을 활용하여 아래 조건에 따라 영업손실보상액을 산정하되, 구체적인 산출근거를 제시하시오. \par
1) 영업허가를 득하고 영업장소가 적법인 경우 \par
2) 영업허가를 득하고 영업장소가 무허가 건축물인 경우 \par
3) 무허가 영업이고 영업장소가 적법인 경우 \par
4) 무허가 영업이고 영업장소가 무허가 건축물인 경우 \par
<자료 1> 사업개요 \par
1) 사업의 종류 : ㅇㅇ 택지개발사업 \par
2) 택지개발사업 예정지구 공람․공고일 : 2006. 04. 05. \par
3) 택지개발사업 개발계획승인 고시일 : 2007. 10. 24. \par
4) 추가 세목고시일 : 2008. 03. 24. \par
5) 협의평가 가격시점 : 2008. 05. 21. \par
6) 재결일 : 2008. 08. 25. \par
7) 현장조사 완료일 : 2008. 09. 21. \par
8) 이의재결시점 : 2008. 10. 05. \par
9) 서울시 강남구, 동작구 및 성남시 수정구와 인접하고 있는 서울시 서초구는 당해 공익사업의 영향으로 지가변동률이 높게 나타나고 있음. \par
10) 당해사업지구의 용도지역이 기존에는 자연녹지(개발제한구역)였으나 공익사업시행에 따른 절차로서 제2종일반주거지역으로 변경되었음. \par
<자료 2> 의뢰물건 내용 \par
1) 토지조서 \par
{\small\centering
\begin{tabular}{|c|c|c|c|c|c|}
\hline
\multirow{2}{*}{\makecell{기호}} & \multirow{2}{*}{\makecell{소재지}} & \multicolumn{2}{|c|}{\makecell{면적(m\textsuperscript{2})}} & \multirow{2}{*}{\makecell{지목}} & \multirow{2}{*}{\makecell{비고}} \\
\cline{3-4}
 &  & \makecell{공부} & \makecell{편입} &  &  \\
\hline
\makecell{1} & \makecell{서초구 신원동 210} & \makecell{450} & \makecell{350} & \makecell{대} &  \\
\hline
\makecell{2} & \makecell{서초구 신원동 221} & \makecell{600} & \makecell{450} & \makecell{대} &  \\
\hline
\makecell{3} & \makecell{서초구 신원동 230} & \makecell{2,000} & \makecell{2,000} & \makecell{임야} &  \\
\hline
\makecell{4} & \makecell{서초구 신원동 240} & \makecell{900} & \makecell{900} & \makecell{전} &  \\
\hline
\end{tabular}\par}

2) 지장물조서 \par
{\small\centering
\begin{tabular}{|c|c|c|c|c|c|}
\hline
\makecell{기호} & \makecell{소재지} & \makecell{물건의 종류} & \makecell{구조․규격} & \makecell{수량} & \makecell{비고} \\
\hline
\makecell{가} & \makecell{신원동 210} & \makecell{주택} & \makecell{시멘트벽돌조 슬래브지붕 단층} & \makecell{50m\textsuperscript{2}} & \makecell{20m\textsuperscript{2} 편입} \\
\hline
\makecell{나} & \makecell{신원동 210} & \makecell{점포} & \makecell{블록조\newline{}스레트지붕 단층} & \makecell{40m\textsuperscript{2}} & \makecell{전부 편입} \\
\hline
\makecell{다} & \makecell{신원동 210} & \makecell{나라안경} & \makecell{-} & \makecell{1식} & \makecell{영업권} \\
\hline
\end{tabular}\par}

<자료 3> 인근지역의 표준지 공시지가 현황 \par
{\small\centering
\begin{tabular}{|c|c|c|c|c|c|c|c|c|c|}
\hline
\multirow{2}{*}{\makecell{기호}} & \multirow{2}{*}{\makecell{소재지}} & \multirow{2}{*}{\makecell{면적\newline{}(m\textsuperscript{2})}} & \multirow{2}{*}{\makecell{지목}} & \multirow{2}{*}{\makecell{이용\newline{}상황}} & \multirow{2}{*}{\makecell{용도\newline{}지역}} & \multirow{2}{*}{\makecell{도로\newline{}교통}} & \multirow{2}{*}{\makecell{형상\newline{}지세}} & \multicolumn{2}{|c|}{\makecell{공시지가(원/m\textsuperscript{2})}} \\
\cline{9-10}
 &  &  &  &  &  &  &  & \makecell{2007년} & \makecell{2008년} \\
\hline
\makecell{A} & \makecell{신원동\newline{}125} & \makecell{300} & \makecell{대} & \makecell{단독} & \makecell{제2종일반} & \makecell{소로한면} & \makecell{세장형 평지} & \makecell{900,000} & \makecell{950,000} \\
\hline
\makecell{B} & \makecell{신원동\newline{}130} & \makecell{900} & \makecell{전} & \makecell{전} & \makecell{제2종일반} & \makecell{세로가} & \makecell{부정형 완경사} & \makecell{600,000} & \makecell{650,000} \\
\hline
\makecell{C} & \makecell{신원동\newline{}산15} & \makecell{3,000} & \makecell{임야} & \makecell{토지임야} & \makecell{제2종일반} & \makecell{맹지} & \makecell{부정형 완경사} & \makecell{250,000} & \makecell{280,000} \\
\hline
\makecell{D} & \makecell{신원동\newline{}233} & \makecell{450} & \makecell{대} & \makecell{단독} & \makecell{개발제한} & \makecell{세로가} & \makecell{가장형 평지} & \makecell{500,000} & \makecell{600,000} \\
\hline
\makecell{E} & \makecell{신원동\newline{}245} & \makecell{450} & \makecell{대} & \makecell{주거나지} & \makecell{개발제한} & \makecell{세로가} & \makecell{세장형 평지} & \makecell{300,000} & \makecell{350,000} \\
\hline
\makecell{F} & \makecell{신원동\newline{}280} & \makecell{1,000} & \makecell{전} & \makecell{전} & \makecell{개발제한} & \makecell{세로\newline{}(불)} & \makecell{부정형 평지} & \makecell{140,000} & \makecell{180,000} \\
\hline
\makecell{G} & \makecell{신원동산100} & \makecell{3,000} & \makecell{임야} & \makecell{토지 임야} & \makecell{개발제한} & \makecell{맹지} & \makecell{부정형 완경사} & \makecell{12,000} & \makecell{15,000} \\
\hline
\end{tabular}\par}

※ 표준지 A, D는 도시계획도로에 20\% 저촉됨. \par
<자료 4> 시점수정 자료 \par
1) 지가변동률 \par
① 서초구, 서울시 평균 용도지역별 지가변동률 (단위 : \%) \par
{\small\centering
\begin{tabular}{|c|c|c|c|c|}
\hline
\multirow{2}{*}{\makecell{구분}} & \multicolumn{2}{|c|}{\makecell{서초구}} & \multicolumn{2}{|c|}{\makecell{서울시 평균}} \\
\cline{2-5}
 & \makecell{주거지역} & \makecell{녹지지역} & \makecell{주거지역} & \makecell{녹지지역} \\
\hline
\makecell{2007. 01. 01. ～ 12. 31.} & \makecell{8.350} & \makecell{10.750} & \makecell{2.350} & \makecell{2.675} \\
\hline
\makecell{2008년 1월} & \makecell{1.100} & \makecell{1.325} & \makecell{0.150} & \makecell{0.235} \\
\hline
\makecell{2008년 2월} & \makecell{1.125} & \makecell{1.355} & \makecell{0.100} & \makecell{0.325} \\
\hline
\makecell{2008년 3월} & \makecell{1.130} & \makecell{1.335} & \makecell{0.125} & \makecell{0.234} \\
\hline
\makecell{2008년 4월} & \makecell{1.145} & \makecell{1.375} & \makecell{0.130} & \makecell{0.235} \\
\hline
\makecell{2008년 5월} & \makecell{1.145} & \makecell{1.375} & \makecell{0.145} & \makecell{0.325} \\
\hline
\makecell{2008년 6월} & \makecell{1.150} & \makecell{1.350} & \makecell{0.125} & \makecell{0.234} \\
\hline
\makecell{2008년 7월} & \makecell{1.100} & \makecell{1.325} & \makecell{0.130} & \makecell{0.225} \\
\hline
\makecell{2008년 8월} & \makecell{1.125} & \makecell{1.355} & \makecell{0.145} & \makecell{0.285} \\
\hline
\end{tabular}\par}

※ 2008년 9월 및 10월의 지가변동률은 미고시된 상태임. \par
② 강남구, 동작구, 성남시 수정구 용도지역별 지가변동률 (단위 : \%) \par
{\small\centering
\begin{tabular}{|c|c|c|c|c|c|c|}
\hline
\multirow{2}{*}{\makecell{구분}} & \multicolumn{2}{|c|}{\makecell{강남구}} & \multicolumn{2}{|c|}{\makecell{동작구}} & \multicolumn{2}{|c|}{\makecell{성남시 수정구}} \\
\cline{2-7}
 & \makecell{주거\newline{}지역} & \makecell{녹지\newline{}지역} & \makecell{주거\newline{}지역} & \makecell{녹지\newline{}지역} & \makecell{주거\newline{}지역} & \makecell{녹지\newline{}지역} \\
\hline
\makecell{2007. 01. 01. ～ 12. 31.} & \makecell{2.350} & \makecell{3.555} & \makecell{2.150} & \makecell{2.750} & \makecell{2.750} & \makecell{2.180} \\
\hline
\makecell{2008. 01. 01. ～ 06. 30.} & \makecell{1.125} & \makecell{2.373} & \makecell{1.145} & \makecell{1.504} & \makecell{1.130} & \makecell{1.565} \\
\hline
\makecell{2008년 7월} & \makecell{0.230} & \makecell{0.335} & \makecell{0.140} & \makecell{0.235} & \makecell{0.120} & \makecell{0.225} \\
\hline
\makecell{2008년 8월} & \makecell{0.245} & \makecell{0.385} & \makecell{0.130} & \makecell{0.275} & \makecell{0.115} & \makecell{0.275} \\
\hline
\end{tabular}\par}

※ 2008년 9월 및 10월의 지가변동률은 미고시된 상태임. \par
2) 생산자물가 상승률 \par
{\small\centering
\begin{tabular}{|c|c|c|c|c|c|c|}
\hline
\makecell{연도} & \makecell{2005. 12.} & \makecell{2006. 01.} & \makecell{2006. 12.} & \makecell{2007. 01.} & \makecell{2007. 12.} & \makecell{2008. 01.} \\
\hline
\makecell{지수} & \makecell{128.8} & \makecell{129.2} & \makecell{130.2} & \makecell{130.5} & \makecell{132.5} & \makecell{132.7} \\
\hline
\makecell{연도} & \makecell{2008. 02.} & \makecell{2008. 03.} & \makecell{2008. 07.} & \makecell{2008. 08.} & \multicolumn{2}{|c|}{\multirow{2}{*}{\makecell{2008년 9월과 10월은 추정.}}} \\
\cline{1-5}
\makecell{지수} & \makecell{132.8} & \makecell{133.0} & \makecell{133.4} & \makecell{133.5} & \multicolumn{2}{|c|}{} \\
\hline
\end{tabular}\par}

<자료 5> 대상물건 조사사항 \par
1) 토지 \par
{\small\centering
\begin{tabularx}{\linewidth}{|c|X|}
\hline
\makecell{구 분} & \multicolumn{1}{c|}{내 용} \\
\hline
\makecell{위치 및 부근상황} & 대상물건은 서초구 신원동 속칭 장수리마을 내에 소재하며 부근은 자연부락내의 단독주택, 농경지, 임야 등으로 형성되어 있음. \\
\hline
\makecell{교통상황} & 대상물건 인근까지 차량접근이 가능하고 인근에 시내버스정류장이 소재하여 일반적인 대중교통사정은 보통임. \\
\hline
\makecell{형태 및 이용상황} & 기호1) 사다리형의 토지로서 남측 인접필지보다 다소 고지이며 주변은 완만한 경사를 이루고 있으며 주상용 건물부지로 이용중임.\newline 기호2) 가장형의 토지로서 인접필지와 등고 평탄하며 주거나지로 이용중임.\newline 기호3) 부정형의 토지로서 서하향의 완경사를 이루고 있으며 전으로 이용중임.\newline 기호4) 세장형의 토지로서 인접필지와 등고 평탄하며 전으로 이용중임. \\
\hline
\makecell{도로상황} & 기호1) 대상물건 남측과 서측으로 각각 노폭 약 8m, 2m 포장도로에 접하고 있음.\newline 기호2) 대상물건 동측으로 노폭 약 2m 포장도로에 접함.\newline 기호3) 지적도상 맹지이나 대상물건 북측으로 약 2m 농로에 접하고 있음.\newline 기호4) 대상물건 동측으로 노폭 약 4m 포장도로에 접함. \\
\hline
\makecell{토지이용계획관계} & 기호1～4) 제2종일반주거지역, 택지개발예정지구, 도시계획도로에 20\%정도 저촉됨. \\
\hline
\makecell{기타 참고사항} & 기호2) 건축허가를 득하였으나 공사착공 전에 사업부지에 편입됨.\newline 기호3) 2001. 03. 05.에 불법으로 형질변경하여 전으로 이용중인 것으로 조사됨.\newline 기호4) 사업구역이 확장되면서 추가로 편입됨. \\
\hline
\end{tabularx}\par}

2) 건물 \par
① 기초자료 \par
{\small\centering
\begin{tabular}{|c|c|c|c|c|}
\hline
\multicolumn{2}{|c|}{\makecell{구분}} & \makecell{기호 가)} & \makecell{기호 나)} & \makecell{비고} \\
\hline
\multicolumn{2}{|c|}{\makecell{사용승인일자}} & \makecell{1989. 10. 25.} & \makecell{1978. 10. 01.} &  \\
\hline
\multicolumn{2}{|c|}{\makecell{건물 내용년수}} & \makecell{45년} & \makecell{40년} &  \\
\hline
\multicolumn{2}{|c|}{\makecell{가격시점 현재\newline{}재조달원가(원/m\textsuperscript{2})}} & \makecell{550,000} & \makecell{450,000} &  \\
\hline
\multirow{6}{*}{\makecell{이전비\newline{}(원)}} & \makecell{해체비} & \makecell{4,000,000} & \makecell{2,000,000} &  \\
\cline{2-5}
 & \makecell{운반비} & \makecell{1,500,000} & \makecell{1,200,000} &  \\
\cline{2-5}
 & \makecell{정지비} & \makecell{1,200,000} & \makecell{1,000,000} &  \\
\cline{2-5}
 & \makecell{재건축비} & \makecell{20,000,000} & \makecell{15,000,000} & \makecell{설비 개량비용\newline{}각 5,000,000원\newline{}포함} \\
\cline{2-5}
 & \makecell{보충자재비} & \makecell{5,000,000} & \makecell{3,000,000} &  \\
\cline{2-5}
 & \makecell{부대비용} & \makecell{5,000,000} & \makecell{3,000,000} &  \\
\hline
\end{tabular}\par}

② 건물 조사내용 \par
- 본 건물의 이전비는 전체 건물을 기준으로 한 것임. \par
- 기호 가) 건물은 기둥이 없는 구조임. \par
- 건물높이는 2m이며 벽면적은 반올림하여 소수점 첫째자리까지 사정함. \par
- 건축 보수비용은 400,000원/m\textsuperscript{2}을 적용함. \par
- 화장실은 편입되어 재설치되어야 하고 위생설비 설치비용은 전체면적을 기준으로 하여 50,000원/m\textsuperscript{2}을 적용함. \par
- 위생설비 이외에는 추가적인 설비공사는 없음. \par
- 건물단가는 천원미만은 절사함. \par
- 건축허가 관련 비용 : 12,000,000원 \par
- 기호 가) 건물 단면도 \par
\begin{center}
\begin{tikzpicture}[x=0.0143cm,y=-0.0143cm,font=\small,line cap=round]
  \fill[gray!45] (410,437) -- (748,406) -- (678,567) -- cycle;
  \draw[very thick] (442,258) -- (748,406) -- (678,567) -- (370,418) -- cycle;
  \draw[very thick,dash pattern=on 6pt off 3pt] (355,441) -- (772,406);
  % dimension lines
  \draw[<->,thin] (452,236) -- (758,384);
  \draw[thin] (442,258) -- (455,230);
  \draw[thin] (748,406) -- (761,378);
  \draw[<->,thin] (417,247) -- (345,407);
  \draw[thin] (442,258) -- (412,244);
  \draw[thin] (370,418) -- (338,404);
  \draw[<->,thin] (773,417) -- (703,578);
  \draw[thin] (748,406) -- (780,420);
  \draw[thin] (678,567) -- (710,581);
  \draw[<->,thin] (398,462) -- (666,592);
  \draw[thin] (410,437) -- (394,470);
  \draw[thin] (678,567) -- (662,600);
  \node[rotate=-25.8,fill=white,inner sep=1pt] at (605,311) {10m};
  \node[rotate=-25.8,fill=white,inner sep=1pt] at (383,322) {5m};
  \node[rotate=-25.8,fill=white,inner sep=1pt] at (738,508) {5m};
  \node[rotate=-25.8,fill=white,inner sep=1pt] at (523,524) {8m};
  \node[rotate=-25.8] at (625,472) {편입부분};
  \node[font=\small\bfseries] at (338,506) {절단선};
  \draw[->,thin] (388,447) -- (360,484);
\end{tikzpicture}
\end{center}
<자료 6> 요인비교 자료 \par
1) 표준지와 보상선례의 지역요인은 동일함. \par
2) 이용상황 \par
{\small\centering
\begin{tabular}{|c|c|c|}
\hline
\makecell{구분} & \makecell{주거용} & \makecell{주상용} \\
\hline
\makecell{주거용} & \makecell{1.00} & \makecell{1.20} \\
\hline
\makecell{주상용} & \makecell{0.83} & \makecell{1.00} \\
\hline
\end{tabular}\par}

3) 도로접면 \par
{\small\centering
\begin{tabular}{|c|c|c|c|c|c|}
\hline
\makecell{구분} & \makecell{소로각지} & \makecell{소로한면} & \makecell{세로가} & \makecell{세로(불)} & \makecell{맹지} \\
\hline
\makecell{소로각지} & \makecell{1.00} & \makecell{0.96} & \makecell{0.86} & \makecell{0.80} & \makecell{0.70} \\
\hline
\makecell{소로한면} & \makecell{1.04} & \makecell{1.00} & \makecell{0.90} & \makecell{0.85} & \makecell{0.75} \\
\hline
\makecell{세로가} & \makecell{1.16} & \makecell{1.05} & \makecell{1.00} & \makecell{0.95} & \makecell{0.83} \\
\hline
\makecell{세로(불)} & \makecell{1.25} & \makecell{1.15} & \makecell{1.11} & \makecell{1.00} & \makecell{0.88} \\
\hline
\makecell{맹지} & \makecell{1.43} & \makecell{1.33} & \makecell{1.20} & \makecell{1.13} & \makecell{1.00} \\
\hline
\end{tabular}\par}

4) 형상 \par
{\small\centering
\begin{tabular}{|c|c|c|c|c|}
\hline
\makecell{구분} & \makecell{가장형} & \makecell{세장형} & \makecell{사다리형} & \makecell{부정형} \\
\hline
\makecell{가장형} & \makecell{1.00} & \makecell{0.97} & \makecell{0.92} & \makecell{0.85} \\
\hline
\makecell{세장형} & \makecell{1.03} & \makecell{1.00} & \makecell{0.95} & \makecell{0.88} \\
\hline
\makecell{사다리형} & \makecell{1.09} & \makecell{1.05} & \makecell{1.00} & \makecell{0.92} \\
\hline
\makecell{부정형} & \makecell{1.18} & \makecell{1.14} & \makecell{1.08} & \makecell{1.00} \\
\hline
\end{tabular}\par}

5) 지세 \par
{\small\centering
\begin{tabular}{|c|c|c|}
\hline
\makecell{구분} & \makecell{평지} & \makecell{완경사} \\
\hline
\makecell{평지} & \makecell{1.00} & \makecell{0.90} \\
\hline
\makecell{완경사} & \makecell{1.11} & \makecell{1.00} \\
\hline
\end{tabular}\par}

6) 도시계획도로 (저촉률 20\%일 경우 전체 토지에 적용하는 비준율) \par
{\small\centering
\begin{tabular}{|c|c|c|}
\hline
\makecell{구분} & \makecell{비저촉} & \makecell{저촉} \\
\hline
\makecell{비저촉} & \makecell{1.00} & \makecell{0.97} \\
\hline
\makecell{저촉} & \makecell{1.03} & \makecell{1.00} \\
\hline
\end{tabular}\par}

<자료 7> 보상평가선례 \par
1) 사업명 : ㅇㅇ 지구 택지개발사업 \par
2) 가격시점 : 2008. 01. 24. \par
3) 소재지 : 강남구 세곡동 424-5번지 \par
4) 용도지역 : 자연녹지지역(개발제한구역) \par
5) 기타요인 보정치 산정을 위한 시점수정치는 1.02100을 적용함. \par
6) 보상단가 등 \par
{\small\centering
\begin{tabular}{|c|c|c|c|}
\hline
\multicolumn{2}{|c|}{\makecell{구분}} & \makecell{424-5번지} & \makecell{500번지} \\
\hline
\multicolumn{2}{|c|}{\makecell{지목}} & \makecell{대} & \makecell{전} \\
\hline
\multicolumn{2}{|c|}{\makecell{단가(원/m\textsuperscript{2})}} & \makecell{700,000} & \makecell{240,000} \\
\hline
\multirow{5}{*}{\makecell{토지\newline{}특성}} & \makecell{이용상황} & \makecell{주상용} & \makecell{전} \\
\cline{2-4}
 & \makecell{도로접면} & \makecell{세로(불)} & \makecell{맹지} \\
\cline{2-4}
 & \makecell{형상} & \makecell{부정형} & \makecell{사다리형} \\
\cline{2-4}
 & \makecell{지세} & \makecell{완경사} & \makecell{완경사} \\
\cline{2-4}
 & \makecell{도시계획도로} & \makecell{저촉} & \makecell{비저촉} \\
\hline
\end{tabular}\par}

<자료 8> 영업보상 관련 자료 \par
1) 대상건물의 임차인은 개인사업자로서 2003. 12. 01.부터 안경점을 운영하여 왔음. \par
2) 영업이익에 관한 자료 \par
① 재무제표에 의한 영업이익 산정 (단위 : 원) \par
{\small\centering
\begin{tabular}{|c|c|c|c|c|}
\hline
\makecell{구분} & \makecell{2004년} & \makecell{2005년} & \makecell{2006년} & \makecell{2007년} \\
\hline
\makecell{매출액} & \makecell{180,000,000} & \makecell{200,000,000} & \makecell{240,000,000} & \makecell{150,000,000} \\
\hline
\makecell{매출원가} & \makecell{87,000,000} & \makecell{95,000,000} & \makecell{113,000,000} & \makecell{65,000,000} \\
\hline
\makecell{판매 및 일반관리비} & \makecell{35,000,000} & \makecell{40,000,000} & \makecell{50,000,000} & \makecell{40,000,000} \\
\hline
\end{tabular}\par}

※ 2007년 매출액은 택지개발사업 개발계획승인이 고시됨으로써 매출액이 감소된 것으로 조사됨. \par
② 부가가치세 과세표준액 기준 매출액 등 \par
{\small\centering
\begin{tabular}{|c|c|c|}
\hline
\makecell{구분} & \makecell{매출액(원)} & \makecell{표준소득률(\%)} \\
\hline
\makecell{2004년} & \makecell{110,000,000} & \makecell{20} \\
\hline
\makecell{2005년} & \makecell{120,000,000} & \makecell{20} \\
\hline
\makecell{2006년} & \makecell{150,000,000} & \makecell{20} \\
\hline
\makecell{2007년} & \makecell{90,000,000} & \makecell{20} \\
\hline
\end{tabular}\par}

③ 인근 동종 유사규모 업종의 영업이익 수준 \par
대상물건을 포함한 인근지역 내 동종 유사규모 업종의 매출액을 탐문조사한 바 연간 220,000,000원 수준이고 매출액 대비 영업이익률은 약 30\%인 것으로 조사되었음. \par
3) 이전 관련자료 \par
① 상품재고액 : 30,000,000원 \par
② 상품운반비 : 3,000,000원 \par
③ 영업시설 등의 이전비 : 2,000,000원 \par
④ 상품의 이전에 따른 감손상당액 : 상품가액의 10\% \par
⑤ 고정적비용 : 임차인은 영업과 관련된 차량에 대한 자동차세 600,000원과 매달 임대료로 500,000원을, 종업원(소득세 원천징수 안함)은 2인으로서 각각 1,200,000원/월을 지급하고 있으며 휴업기간 중에는 1인만 필요함. \par
⑥ 이전광고비 및 개업비 등 부대비용 : 2,000,000원 \par
4) 기타자료 \par
① 제조부문 보통인부 노임단가 : 50,000원/일 \par
② 도시근로자 월평균 가계지출비 \par
{\small\centering
\begin{tabular}{|c|c|}
\hline
\makecell{구 분} & \makecell{월평균 가계지출비} \\
\hline
\makecell{2인} & \makecell{2,500,000} \\
\hline
\makecell{3인} & \makecell{3,000,000} \\
\hline
\makecell{4인} & \makecell{3,500,000} \\
\hline
\makecell{5인} & \makecell{4,000,000} \\
\hline
\makecell{6인} & \makecell{4,500,000} \\
\hline
\end{tabular}\par}

③ 영업이익은 만원단위에서 반올림하여 사정함. \par
\end{longquestion}
